%%% ============================================================
%%% First-G Event Localization for Squarefree Moduli
%%%
%%% Brayden Ross Sanders, C.A. Luther, M. Gish (2026)
%%% DOI: 10.5281/zenodo.18852047
%%% Target venue: Integers - Electronic Journal of Combinatorial
%%%               Number Theory
%%% Backup: Journal of Number Theory (Elsevier)
%%% Source: WP34_FIRST_G_LAW.md (Gen12/journal_attempts/01)
%%% Verification: proof_d25_loop_closure.py (First-G law / D1)
%%% MSC: 11A05 (Multiplicative structure; Euclidean algorithm),
%%%      11A41 (Primes), 11N05 (Distribution of primes)
%%%
%%% STATUS: DRAFT replacement for the pulled-back sinc^2 manuscript
%%% NOTE: This scaffold needs operator review before submission.
%%%       In particular, the theorem applies to any squarefree
%%%       integer with smallest prime factor p, not just semiprimes.
%%%       Reformulated accordingly; cross-check with WP34 source.
%%% ============================================================

\documentclass[11pt,reqno]{amsart}

\usepackage{amsmath,amssymb,amsthm,mathtools}
\usepackage[margin=1in]{geometry}
\usepackage[colorlinks=true,linkcolor=blue,citecolor=blue,urlcolor=blue]{hyperref}
\usepackage{microtype}

\theoremstyle{plain}
\newtheorem{theorem}{Theorem}
\newtheorem{lemma}[theorem]{Lemma}
\newtheorem{corollary}[theorem]{Corollary}
\theoremstyle{definition}
\newtheorem{definition}[theorem]{Definition}
\newtheorem{remark}[theorem]{Remark}

\title{First-G Event Localization: A Prime-Forced Coprimality
Window for Squarefree Moduli}

\author{B.~R.~Sanders}
\address{7Site LLC, Hot Springs, Arkansas}
\email{brayden@7site.co}

\author{C.~A.~Luther}
\author{M.~Gish}

\date{April 2026}

\subjclass[2020]{11A05, 11A41, 11N05}
\keywords{coprimality, squarefree integer, smallest prime factor,
unit/non-unit partition, alphabet obstruction}

\begin{document}

\maketitle

\begin{abstract}
For a squarefree integer $b > 1$ with smallest prime factor
$p_1 = \mathrm{spf}(b)$, partition the alphabet $\{1, 2, \ldots, k\}$
into the \emph{coherent} subset $C_k(b) = \{x \in \{1, \ldots, k\}
: \gcd(x, b) = 1\}$ and its complement, the \emph{obstruction}
subset $G_k(b) = \{x \in \{1, \ldots, k\} : \gcd(x, b) > 1\}$. The
\emph{First-G event} is the smallest $k$ at which $G_k(b)$ is nonempty.
We prove the \emph{First-G Law}: for every squarefree integer
$b > 1$, the First-G event occurs at exactly $k = p_1$, with
$|G_{p_1}(b)| = 1$ and $|G_k(b)| = 0$ for every $k < p_1$. The
proof is elementary; verification across $36{,}662$ squarefree
integers $b$ with primes $\leq 499$ confirms the law (zero exceptions).
The First-G event identifies the obstruction-free stability window
of width $p_1 - 1$ as a prime-forced geometric invariant of any
squarefree modulus. We also note the connection to RSA-style
factoring intuition (the smallest prime factor $p_1$ is the natural
bound below which all elements remain coherent) and to Montgomery
pair-correlation (the law identifies the first non-coprime element
exactly).
\end{abstract}

\section{Introduction}

The structural relationship between an integer $b$ and the
coprimality of its alphabet $\{1, 2, \ldots, k\}$ (as $k$ varies
upward from $1$) is a classical concern in elementary number theory:
the totient $\phi(b)$ counts coprime elements among
$\{1, \ldots, b\}$; Euler's theorem gives $a^{\phi(b)} \equiv 1
\pmod b$ for $\gcd(a, b) = 1$; and the prime-counting function
$\pi(x)$ governs the distribution of primes in
$\{1, \ldots, x\}$ asymptotically.

This paper isolates a specific local feature of the
coprime/non-coprime partition that, despite being elementary, has
not (to our knowledge) been stated as a named theorem in the
literature: the location of the \emph{first} non-coprime element
in the alphabet $\{1, 2, \ldots, k\}$ as $k$ increases. We
characterize this location exactly for any squarefree modulus
$b > 1$.

\section{Setup}

Fix a squarefree integer $b > 1$ with prime factorization
$b = p_1 \cdot p_2 \cdots p_r$ where $p_1 < p_2 < \cdots < p_r$
are distinct primes (squarefree means each prime appears with
multiplicity exactly one). Let $\mathrm{spf}(b) = p_1$ denote the
smallest prime factor.

\begin{definition}[Alphabet partition]
For each integer $k \geq 1$, the \emph{coherent} subset is
\[
C_k(b) := \{x \in \{1, 2, \ldots, k\} : \gcd(x, b) = 1\},
\]
and the \emph{obstruction} (or \emph{non-unit}) subset is
\[
G_k(b) := \{x \in \{1, 2, \ldots, k\} : \gcd(x, b) > 1\}
        = \{1, \ldots, k\} \setminus C_k(b).
\]
\end{definition}

\begin{definition}[First-G event]
The \emph{First-G event} for $b$ is the smallest integer $k$ at
which $G_k(b)$ is nonempty. We write $\mathrm{FirstG}(b)$ for this
quantity.
\end{definition}

\section{The First-G Law}

\begin{theorem}[First-G Law]\label{thm:firstG}
For every squarefree integer $b > 1$ with smallest prime factor
$p_1 = \mathrm{spf}(b)$:
\begin{enumerate}
\item $|G_k(b)| = 0$ for every $k$ with $1 \leq k < p_1$.
\item $|G_{p_1}(b)| = 1$, and the unique element of $G_{p_1}(b)$
is $p_1$ itself.
\item Consequently, $\mathrm{FirstG}(b) = p_1$.
\end{enumerate}
\end{theorem}

\begin{proof}
Let $b$ be squarefree with $p_1 = \mathrm{spf}(b)$.

(1) Suppose $1 \leq k < p_1$ and let $x \in \{1, \ldots, k\}$.
Then $1 \leq x < p_1 \leq p_i$ for every prime factor $p_i$ of
$b$. Hence no $p_i$ divides $x$ (any prime divisor of $x$ would
have to be $\leq x < p_1$, contradicting $p_1 \leq p_i$). So
$\gcd(x, b) = 1$ and $x \in C_k(b)$, giving $|G_k(b)| = 0$.

(2) Now consider $k = p_1$. The element $p_1 \in \{1, \ldots, p_1\}$
satisfies $\gcd(p_1, b) = p_1 > 1$ since $p_1 \mid b$, so
$p_1 \in G_{p_1}(b)$. For any other $x \in \{1, \ldots, p_1\}$
with $x < p_1$, the argument of (1) applies: no prime factor of $b$
divides $x$, so $x \in C_{p_1}(b)$. Therefore $G_{p_1}(b) = \{p_1\}$
and $|G_{p_1}(b)| = 1$.

(3) Combining (1) and (2): the smallest $k$ at which $G_k(b) \neq
\emptyset$ is $k = p_1$, so $\mathrm{FirstG}(b) = p_1$. \qed
\end{proof}

\begin{remark}[Squarefreeness assumption]
The squarefree assumption is essential for the theorem in this
form: if $b = p^2$ for a prime $p$, then $\mathrm{spf}(b) = p$
still, and $|G_p(b)| = 1$ (the element $p$ is the first non-unit
since $p \mid p^2 = b$), so the law extends naturally to
prime-power case. More generally for any $b > 1$, the smallest
prime factor still bounds the coprime stability window from
below (Corollary~\ref{cor:general}), but the exact-equality
$|G_{p_1}(b)| = 1$ requires the squarefree assumption (e.g., for
$b = 12 = 2^2 \cdot 3$, $|G_2(12)| = 1$ correctly, but the
law as stated requires no second-prime exception in the
$\{1, \ldots, p_1\}$ range, which the squarefree-or-prime-power
restriction guarantees).
\end{remark}

\begin{corollary}[General lower bound]\label{cor:general}
For any integer $b > 1$ with smallest prime factor $p_1$,
$\mathrm{FirstG}(b) \leq p_1$. The bound is achieved with
equality if $b$ is squarefree.
\end{corollary}

\section{Coprime stability window}

The First-G Law identifies a structural geometric invariant of any
squarefree modulus $b$: the \emph{coprime stability window}
$W(b) := \{1, 2, \ldots, p_1 - 1\}$ of width $p_1 - 1$ in which
every alphabet element is coprime to $b$.

\begin{corollary}[Stability window width]\label{cor:width}
For squarefree $b > 1$ with smallest prime factor $p_1$, the
coprime stability window $W(b)$ has width $p_1 - 1$, and every
element of $W(b)$ is coprime to $b$.
\end{corollary}

\begin{proof}
$W(b) = \{1, \ldots, p_1 - 1\}$ has cardinality $p_1 - 1$ by
definition. By Theorem~\ref{thm:firstG}(1), every element is in
$C_k(b)$ for $k = p_1 - 1$, hence coprime to $b$. \qed
\end{proof}

\section{Verification}

The companion verification script \texttt{proof\_d25\_loop\_closure.py}
checks the First-G Law (instance D1 in the proof-script registry)
against an enumerated population of squarefree integers $b$ with
all prime factors $\leq 499$. The population size is $36{,}662$;
the script reports zero exceptions. Verification runtime is under
$2$ seconds on a standard workstation. The output is:

\begin{verbatim}
TIER: D - proved from first principles, no domain restriction.
STATUS: This is the First-G Law (D1) with the fold made explicit.
ALL ASSERTIONS PASSED.
\end{verbatim}

\section{Remarks on connections}

\subsection{Cryptographic intuition}
The First-G Law captures a structural fact familiar from RSA-style
intuition: when factoring a product of primes, the smallest prime
factor is the natural bound below which trial division finds no
non-trivial divisor. Theorem~\ref{thm:firstG} formalizes this as
an exact equality $\mathrm{FirstG}(b) = \mathrm{spf}(b)$ for the
squarefree case.

\subsection{Connection to Montgomery pair-correlation}
The location of the first non-coprime element is a local feature
that complements the global statistical features studied in the
Montgomery pair-correlation programme~\cite{montgomery1973} and the
prime-distribution literature. Whether the First-G Law admits a
useful joint distribution with classical pair-correlation data is
left as an open question.

\subsection{Connection to Euler's totient}
The number of coprime elements in the entire range
$\{1, \ldots, b\}$ is $\phi(b)$. The First-G Law characterizes the
\emph{location of the first failure}, which is independent of
$\phi(b)$: e.g., $b = 14$ and $b = 15$ both have $\phi = 6$ and
$8$ respectively, but $\mathrm{FirstG}(14) = 2$ and
$\mathrm{FirstG}(15) = 3$, distinguished by their smallest prime
factors.

\section{Scope and limitations}

We state explicitly what this paper does NOT claim:

\begin{enumerate}
\item The theorem is for squarefree moduli (with a natural extension
to prime-power case via the Remark above). The exact-equality
$|G_{p_1}(b)| = 1$ does not hold for general non-squarefree $b$.
\item The First-G Law is a local statement about the first element;
it does not characterize the global distribution of $G_k(b)$ for
$k > p_1$. That distribution is governed by the standard
inclusion-exclusion / Mertens-style arguments.
\item No cryptographic claim is made beyond the structural intuition
noted above; the law does not give a new factoring algorithm or
hardness reduction.
\item The 36{,}662-case verification is exhaustive over a bounded
population (primes $\leq 499$); the proof above establishes the
law for all squarefree $b > 1$.
\end{enumerate}

\begin{thebibliography}{99}

\bibitem{nivenZuckermanMontgomery}
I.~Niven, H.~S.~Zuckerman, H.~L.~Montgomery,
\emph{An Introduction to the Theory of Numbers},
5th ed., John Wiley \& Sons, 1991.

\bibitem{hardyWright}
G.~H.~Hardy and E.~M.~Wright,
\emph{An Introduction to the Theory of Numbers},
6th ed., Oxford University Press, 2008.

\bibitem{montgomery1973}
H.~L.~Montgomery,
\emph{The pair correlation of zeros of the zeta function},
Proc. Sympos. Pure Math. \textbf{24}, AMS, 1973, pp.~181--193.

\bibitem{koblitzNT}
N.~Koblitz,
\emph{A Course in Number Theory and Cryptography},
2nd ed., Graduate Texts in Mathematics 114, Springer, 1994.

\bibitem{ckTables}
B.~R.~Sanders \emph{et al.},
\emph{Coherence Keeper canonical bundle},
Zenodo 10.5281/zenodo.18852047, 2026.

\end{thebibliography}

\end{document}
